\documentclass[11pt,a4paper]{article}

\usepackage[utf8]{inputenc}
\usepackage[T1]{fontenc}
\usepackage{amsmath,amssymb,amsfonts}
\usepackage{amsthm}
\usepackage{mathtools}
\usepackage{bm}

\theoremstyle{plain}
\newtheorem{proposition}{Proposition}
\newtheorem{corollary}{Corollary}
\theoremstyle{definition}
\newtheorem{assumption}{Assumption}
\newtheorem{definition}{Definition}
\usepackage{geometry}
\usepackage{booktabs}
\usepackage{longtable}
\usepackage{enumitem}
\usepackage[hidelinks]{hyperref}
\geometry{margin=2.6cm}

\newcommand{\R}{\mathbb{R}}
\newcommand{\Rpos}{\mathbb{R}_{\geq 0}}
\newcommand{\Z}{\mathbb{Z}}
\newcommand{\Zpos}{\mathbb{Z}_{\geq 0}}
\newcommand{\E}{\mathbb{E}}
\newcommand{\pa}{\mathrm{pa}}
\newcommand{\diff}{\mathrm{d}}

\title{A Multi-Stage Stochastic Capacity-Expansion Model for the Sizing of
Off-Grid Solar Microgrids under Demand, Climate and Economic Uncertainty}
\author{}
\date{}

\begin{document}
\maketitle

\begin{abstract}
This document develops the mathematical formulation of a capacity-expansion model
for the optimal sizing of off-grid solar microgrids serving rural communities. The
system couples a photovoltaic array, a diesel generator, a
battery and a hybrid inverter to a community load. Investment and operation are
exposed to four families of uncertainty---future demand, the renewable resource,
economic conditions and policy---which are represented by a discrete scenario tree.
The model is posed as an adaptive multi-stage stochastic program in which capacity is
expanded incrementally at successive milestone years and operation is dispatched by a
rule-based generation-balance controller rather than by cost minimisation. We give a
self-contained statement of the sets, parameters, decision variables, constraints and
objective, we show how the rule-based controller is reproduced exactly through a small
number of mixed-integer constraints, and we describe the construction and reduction of
the scenario tree. The resulting deterministic-equivalent problem is a mixed-integer
linear program of moderate size. Beyond the formulation, we establish that the
cost-optimal dispatch problem is a valid lower bound on the operating cost of the fixed
rule-based controller, which lets us certify the global optimality of the sizing for the
controller that is actually deployed and quantify the ``price of the heuristic
dispatch''---a certified gap that, to our knowledge, no existing off-grid sizing method
provides.
\end{abstract}

\tableofcontents

\section{Context}
\label{sec:intro}

\subsection{Background and problem}

Stand-alone (off-grid) solar microgrids---coupling a photovoltaic array, a diesel
generator, battery storage and a hybrid inverter to a community load---are a leading
option for electrifying rural communities beyond the reach of the main grid. Their
economics are dominated by the initial and staged investment in generation and storage
capacity: the sizing decision fixes most of the life-cycle cost and the quality of
service delivered. Sizing is difficult for two reasons. First, the conditions that
determine the value of each asset are deeply uncertain over the multi-decade project
horizon: demand grows as households connect and adopt appliances and productive uses, the
solar resource shifts with the climate, and fuel and equipment prices follow uncertain
trajectories. Second, the cost of a given fleet depends on how it is \emph{operated}, and
in the field operation is governed not by an optimiser but by a fixed, low-cost
rule-based controller embedded in the power electronics. A credible sizing method must
therefore contend with both the uncertainty and the operating policy that will actually
run the plant.

\subsection{Related work}
\label{sec:relwork}

The literature on off-grid microgrid sizing separates into two families distinguished by
how they treat operation and whether they can certify optimality.

The first family relies on \emph{rule-based} dispatch and searches the capacity space
heuristically. The de-facto industry standard, HOMER and HOMER Pro, dispatches the plant
with fixed strategies---load following, cycle charging, generator order, combined and
predictive dispatch---and locates capacities either by enumerating a user-defined grid or
by a proprietary derivative-free optimiser that, in the vendor's own words, can risk
converging to a local optimum. The widely used iHOGA/MHOGA tool of Dufo-L\'opez and
Bernal-Agust\'in optimises both the sizing and the control strategy with genetic
algorithms, returning ``suitable'' rather than certified solutions. Research tools in the
same vein, including the uGrid/uGridNet generation-balance controller adopted in the
present work, likewise pair a rule-based dispatcher with metaheuristic or grid search.
Across this family the operating model is realistic---it is the controller that will run
the plant---but the reported sizing carries \emph{no optimality certificate, gap or
bound}.

The second family uses \emph{exact mathematical programming}. Mixed-integer linear and
non-linear programs co-optimise sizing and dispatch and return a globally optimal (or
bounded) solution for the formulation solved: representative examples are the MILP
planning model of Wu \emph{et al.} for a stand-alone photovoltaic/diesel/battery microgrid
(IEEE Transactions on Smart Grid 12(5), 2021), the two-stage stochastic mixed-integer
programme of Akulker \emph{et al.} (2025), the multi-year MicroGridsPy capacity-expansion
framework of Stevanato \emph{et al.} (Energy for Sustainable Development, 2020), and the
multi-stage stochastic integer template of Ahmed, King and Parija (Journal of Global
Optimization 26, 2003). Decomposition methods can additionally report an optimality gap;
the progressive-hedging scheme in NREL's REopt (Zolan/Wales \emph{et al.}, Optimization
and Engineering 26, 2025) is representative. In every case, however, operation is modelled
as \emph{cost-optimal} economic dispatch: the second-stage decisions are chosen by the
optimiser to minimise short-run cost.

A small number of hybrids sit between the two families but do not close the gap below.
The inner--outer co-optimisation of Lyu \emph{et al.} (Applied Energy 363, 2024) places a
metaheuristic in the outer sizing loop and a cost-optimal energy-management optimiser in
the inner loop, and explicitly deems it ``risky to determine the optimal sizes by
rule-based strategies''. The recent work of Nespoli and Medici (arXiv:2511.21619, 2025)
tunes a rule-based battery controller jointly with its sizing using differential
evolution, but provides no certificate and explicitly dismisses a nested
master/inner-optimisation as too costly; it also addresses grid-connected peak shaving
rather than off-grid supply.

\subsection{Research gap}
\label{sec:gap}

Two observations follow. The methods that use the \emph{deployed} rule-based controller
cannot certify that the chosen sizing is optimal, and typically explore the capacity space
by enumeration or metaheuristics without any bound. The methods that \emph{can} certify
optimality do so for a cost-optimal dispatch that will not run the plant; their optimal
sizing is optimal for the wrong operating model, and the discrepancy---the extra
life-cycle cost incurred because the field controller is heuristic rather than
optimal---is neither certified nor quantified. Consequently:
\begin{quote}
\emph{No published method certifies the global optimality of an off-grid microgrid sizing
for a fixed, non-optimising rule-based dispatch controller, nor quantifies the cost of
using that controller instead of optimal dispatch at the design stage.}
\end{quote}
This is the gap the present work addresses.

\subsection{Hypotheses}
\label{sec:hypotheses}

The approach rests on three hypotheses.
\begin{description}[leftmargin=1.6em,style=nextline]
  \item[H1 (bounding).] For any fixed capacity and scenario, the operating cost under
  cost-optimal dispatch is no greater than the operating cost under the rule-based
  controller, because the controller's trajectory is one feasible dispatch among those the
  cost-optimal problem minimises over. The cost-optimal sizing problem is therefore a valid
  \emph{global lower bound} on the rule-based sizing problem.
  \item[H2 (certification).] Because a computable lower bound (the cost-optimal relaxation)
  can be paired with a computable upper bound (a forward simulation of the actual
  controller), the global optimum of the sizing \emph{for the rule-based controller} can be
  certified by branch-and-bound over the discrete capacity lattice---without ever encoding
  the controller as constraints or approximating its behaviour.
  \item[H3 (price of the heuristic).] The difference between the rule-based and the
  cost-optimal sizing optima is a well-defined, non-negative and computable quantity---the
  \emph{price of the heuristic dispatch}---that measures the value of dispatch optimality at
  the design stage.
\end{description}
Hypothesis~H1 requires only that the controller respect the same physical constraints as
the cost-optimal model; this is stated precisely as Assumption~\ref{ass:feas} and verified
for the generation-balance controller in Section~\ref{sec:certify}.

\subsection{Methodology and contributions}
\label{sec:method}

To fill the gap the document develops, in order:
\begin{enumerate}[leftmargin=1.8em,itemsep=1pt]
  \item an \emph{adaptive multi-stage stochastic capacity-expansion} model on a scenario
  tree (Sections~\ref{sec:framework}--\ref{sec:objective}) that represents the four
  families of uncertainty---demand, renewable resource, economics and policy---and expands
  capacity incrementally at milestone years under non-anticipativity;
  \item two operating models sharing one feasible set: the rule-based generation-balance
  controller, reproduced exactly as a binary-light mixed-integer encoding
  (Section~\ref{sec:dispatch}), and its cost-optimal relaxation, which serves as the
  lower-bounding oracle;
  \item the construction and reduction of the scenario tree and the time domain
  (Section~\ref{sec:scenarios}), yielding a tractable deterministic-equivalent mixed-integer
  linear program (Section~\ref{sec:solution});
  \item the central methodological contribution (Section~\ref{sec:certify}): the
  dispatch-relaxation bound of~H1 (Proposition~\ref{prop:bound}), the
  \emph{branch-and-simulate} algorithm that certifies the rule-based sizing optimum
  of~H2 (Corollary~\ref{cor:cert}), and the price of the heuristic dispatch of~H3
  (Definition~\ref{def:price}).
\end{enumerate}
Taken together these bridge the two families of Section~\ref{sec:relwork}: they retain the
\emph{realistic} rule-based controller of the first while delivering the \emph{certified}
optimality of the second, and they make the cost of the heuristic itself a reported
result.

\section{System and modelling framework}
\label{sec:framework}

We consider a stand-alone (off-grid) microgrid whose generation and storage assets
must be sized, and subsequently expanded, so as to supply a growing rural community at
least cost while maintaining an acceptable quality of service. Four technologies are
available. A photovoltaic array provides non-dispatchable renewable generation; a
diesel generator provides dispatchable back-up generation; a
battery provides storage; and a hybrid inverter performs the alternating- and
direct-current conversions that couple these assets to the load. The community load is
the only demand, and there is no connection to a larger grid.

The planner faces uncertainty that resolves only gradually over the lifetime of the
project. The community's demand grows as households connect, acquire appliances and
change their usage habits as well as productive usages ; the solar resource evolves with the climate; fuel prices
and equipment costs follow uncertain trajectories; and policy may impose a minimum
renewable penetration. Because the planner can observe how these conditions unfold and
react by expanding capacity at later milestone years, the appropriate decision-making
paradigm is \emph{adaptive multi-stage stochastic programming}. The uncertain future
is discretised into a scenario tree (Section~\ref{sec:scenarios}); investment decisions
are attached to the nodes of this tree and are therefore allowed to depend on, but only
on, the information revealed up to the stage at which they are taken. This information
restriction---non-anticipativity---is enforced structurally by the tree indexing
rather than by explicit coupling constraints.

A distinctive feature of the present model is its treatment of operation. In the
classical capacity-expansion formulation the second-stage operating problem is itself an
optimisation that dispatches each asset to minimise short-run cost. Here we instead
adopt the rule-based ``generation-balance control'' used in off-grid practice, in which
the dispatch follows a fixed priority order---solar first, then the battery, and the
diesel generator only as a last resort---together with a look-ahead rule that reserves
enough stored energy to carry the community through the night. We show in
Section~\ref{sec:dispatch} that this controller can be reproduced within a mixed-integer
linear program using a single binary variable per time step, so that the sizing problem
remains a single mathematical program whose solution is consistent with the operating
policy that will actually be used in the field.

Throughout, the planning horizon is partitioned into stages indexed by milestone years.
Within each tree node a full year of operation is represented by a small set of
weighted representative days, which keeps the operating sub-problem tractable while
preserving the diurnal and seasonal structure that drives the value of storage. The
remainder of the document states the model formally.

\section{Sets and indices}
\label{sec:sets}

The scenario tree is the backbone of the formulation. Its nodes carry both the
investment decisions and the operating sub-problems, and the economics are aggregated
over it.

\begin{description}[leftmargin=2.2cm,style=nextline]
  \item[$\mathcal{N}$] Set of nodes of the scenario tree, indexed by $n$. Each node
  represents one realisation of the uncertain conditions up to its stage. The unique
  root node is denoted $n_0$.
  \item[$\mathcal{G}$] Set of technologies, $\mathcal{G}=\{\mathrm{pv},\,
  \mathrm{batt},\,\mathrm{ge},\,\mathrm{inv}\}$, denoting respectively the photovoltaic
  array, the battery, the diesel generator and the inverter. 
  \item[$\mathcal{T}_n$] Set of representative days associated with node $n$, indexed by
  $t$. Each representative day stands for a cluster of similar days within the operating
  year that the node represents.
  \item[$\mathcal{H}$] Set of intra-day time steps, indexed by $h\in\{0,\dots,23\}$
  (hourly resolution). Operation is defined on the triples $(n,t,h)$.
  \item[$\mathcal{S}_{\mathrm{ge}}$] Discrete catalogue of commercially available
  generator ratings, indexed by $s$. The generator is a \emph{single} unit chosen from
  this catalogue and replaced when it is upgraded, whereas the photovoltaic array, the
  battery and the inverter are \emph{modular} and expanded in identical units.
\end{description}

Three structural maps are defined on the tree. For a non-root node $n$, $\pa(n)$ denotes
its parent; $\mathrm{stage}(n)\in\{0,1,\dots,S\}$ denotes the stage (milestone year) at
which $n$ sits; and $\mathcal{L}\subset\mathcal{N}$ denotes the set of leaf nodes, which
correspond to complete scenarios. The path probability $P_n$ of a node is the
probability that the realised scenario passes through $n$; by construction
$\sum_{\ell\in\mathcal{L}}P_\ell=1$.

\section{Parameters}
\label{sec:parameters}

Parameters are grouped into those describing the operating environment of a node, the
technical behaviour of the assets, and the economics. Time-varying environmental
parameters are indexed by $(n,t,h)$ because each tree node carries its own realisation
of demand and resource.

\paragraph{Notation convention.} Throughout the formulation, \emph{uppercase} symbols
denote parameters---data fixed before optimisation---and \emph{lowercase} symbols denote
decision variables, so that the two can be told apart at a glance in every equation. The
only exceptions are the policy hyper-parameters $\theta\in\Theta$ of
Section~\ref{sec:refine}, which follow the conventional control-theoretic notation, and a
handful of named set-valued maps ($\pa$, $\mathrm{stage}$, $\mathrm{yr}$).

\paragraph{Demand and resource.}
\begin{description}[leftmargin=3.0cm,style=nextline]
  \item[$D_{n,t,h}$] Community electrical demand in time step $(n,t,h)$ [kW].
  \item[$Y_{n,t,h}$] Specific photovoltaic yield, i.e.\ the power produced per unit of
  installed array capacity [kW\,kW$^{-1}$], obtained from the node's irradiance scenario.
  \item[$W_{n,t}$] Weight of representative day $t$ in node $n$ [days\,yr$^{-1}$], equal
  to the number of actual days the representative day stands for; $\sum_{t\in
  \mathcal{T}_n} W_{n,t}=365$.
\end{description}

\paragraph{Technical parameters.}
\begin{description}[leftmargin=3.0cm,style=nextline]
  \item[$H^{\mathrm{c}},\ H^{\mathrm{d}}$] One-way charging and discharging
  efficiencies of the battery.
  \item[$A_{n,t,h}$] Fractional self-discharge of the battery over one time step,
  dependent on ambient temperature.
  \item[{$F^{e}_{n,t,h}\in(0,1]$}] Temperature-dependent usable-capacity factor of the
  battery, which scales the upper energy limit.
  \item[$\underline{E},\ \overline{E}$] Lower and upper state-of-charge fractions
  ($0.05$ and $0.95$), expressing the protective trip limits of the battery management
  system.
  \item[$R_{n,t,h}$] Look-ahead night-reserve energy [kWh]: a forecast of the demand
  remaining until the next sunrise, used by the generation-balance controller to decide
  when the battery may serve the load (Section~\ref{sec:dispatch}).
  \item[$F^{\mathrm{ge}}$] Minimum stable loading of the generator as a fraction of
  its rating (taken as $0.30$).
  \item[$F_0,\ F_1$] Intercept and slope of the linearised generator fuel
  characteristic [L\,h$^{-1}$ and L\,kWh$^{-1}$], so that the fuel flow is
  $F_0+F_1 p^{\mathrm{ge}}$ when the generator is on.
  \item[$K^{\mathrm{ge}}_s$] Rated power of generator catalogue option $s$ [kW].
  \item[$U^{\mathrm{pv}},\ U^{\mathrm{batt}},\ U^{\mathrm{inv}}$] Unit increments of the
  modular technologies: the rated power of a single photovoltaic panel [kW], the usable
  energy of a single battery module [kWh] and the rated power of a single inverter module
  [kW].
\end{description}

\paragraph{Economic parameters.}
\begin{description}[leftmargin=3.0cm,style=nextline]
  \item[$C^{\mathrm{inv}}_{g,n}$] Investment (capital) cost of technology $g$ if
  installed at node $n$ [\$\,kW$^{-1}$ or \$\,kWh$^{-1}$], following the node's cost
  trajectory.
  \item[$O_{g}$] Annual operation-and-maintenance cost of technology $g$, expressed as a
  fraction of its capital cost.
  \item[$C^{\mathrm{fuel}}_n$] Diesel price at node $n$ [\$\,L$^{-1}$].
  \item[$C^{\mathrm{deg}}$] Battery throughput-degradation cost [\$\,kWh$^{-1}$ of energy
  discharged], equal to the replacement cost divided by the lifetime energy throughput.
  \item[$V^{\mathrm{lol}}$] Value of lost load [\$\,kWh$^{-1}$], the penalty applied to unmet demand.
  \item[$I$] Annual discount rate.
  \item[$B_n$] Discount factor that brings costs incurred at node $n$ to present
  value, $B_n=(1+I)^{-\mathrm{yr}(n)}$ where $\mathrm{yr}(n)$ is the calendar year of $\mathrm{stage}(n)$.
  \item[$\Delta_n$] Number of operating years that node $n$ represents, i.e.\ the span
  between its stage and the next milestone.
  \item[$C^{g}_0$] Existing installed capacity of technology $g$ already on site at
  commissioning (brownfield initial condition); zero for a greenfield project. 
  \item[$V^{\mathrm{ge}}_s$] Resale (or redeployment/transfer) value recovered when an
  incumbent generator of size $s$ is retired and replaced, net of removal and depreciated
  by the unit's vintage.
\end{description}

\section{Decision variables}
\label{sec:variables}

The decision variables divide naturally into an investment layer, which is attached to
the tree nodes and constitutes the adaptive ``here-and-now'' decisions, and an operating
layer, which is defined on the representative time steps of each node.

\paragraph{Investment layer.} For each node $n$:
\begin{align}
  b^{\mathrm{pv}}_n,\ b^{\mathrm{batt}}_n,\ b^{\mathrm{inv}}_n &\in \Zpos
    && \text{modules of PV / battery / inverter added at $n$,}\\
  z^{\mathrm{ge}}_{n,s} &\in \{0,1\}
    && \text{indicator that size $s$ is the single active generator at $n$,}\\
  i^{\mathrm{ge}}_{n,s},\ r^{\mathrm{ge}}_{n,s} &\in [0,1]
    && \text{install / retire of a size-$s$ generator at $n$,}\\
  \mathrm{cap}^{g}_n &\in \Rpos
    && \text{installed capacity of technology $g$ available at $n$.}
\end{align}
The photovoltaic array, battery and inverter are modular: their capacity accumulates
through the integer counts $b^{\mathrm{pv}},b^{\mathrm{batt}},b^{\mathrm{inv}}$. The
generator is a single unit: exactly one catalogue size is active at each node,
$\sum_{s}z^{\mathrm{ge}}_{n,s}=1$, and an upgrade installs a larger unit while retiring the
incumbent (Section~\ref{sec:brownfield}). The transition variables
$i^{\mathrm{ge}}_{n,s}$ and $r^{\mathrm{ge}}_{n,s}$ record, respectively, a size
newly installed at $n$ and a size retired at $n$.

\paragraph{Operating layer.} For each time step $(n,t,h)$:
\begin{align}
  p^{\mathrm{ge}}_{n,t,h} &\in \Rpos && \text{generator power output [kW],}\\
  p^{\mathrm{c}}_{n,t,h},\ p^{\mathrm{d}}_{n,t,h} &\in \Rpos
    && \text{battery charging / discharging power [kW],}\\
  e_{n,t,h} &\in \Rpos && \text{battery energy content (state of charge) [kWh],}\\
  q_{n,t,h} &\in \Rpos && \text{curtailed photovoltaic power [kW],}\\
  \ell_{n,t,h} &\in \Rpos && \text{unmet demand (loss of load) [kW],}\\
  y_{n,t,h} &\in \{0,1\} && \text{generator commitment (on/off).}
\end{align}
The binary commitment $y_{n,t,h}$ is the only integer operating variable, and there is
exactly one such variable per time step.

\section{Constraints}
\label{sec:constraints}

\subsection{Capacity accumulation and non-anticipativity}
\label{sec:capacity}

The modular technologies accumulate: the capacity available at a node equals the capacity
inherited from its parent plus the modules added at the node,
\begin{align}
  \mathrm{cap}^{\mathrm{pv}}_n   &= \mathrm{cap}^{\mathrm{pv}}_{\pa(n)}
      + U^{\mathrm{pv}}\, b^{\mathrm{pv}}_n, \label{eq:cap-pv}\\
  \mathrm{cap}^{\mathrm{batt}}_n &= \mathrm{cap}^{\mathrm{batt}}_{\pa(n)}
      + U^{\mathrm{batt}}\, b^{\mathrm{batt}}_n, \label{eq:cap-batt}\\
  \mathrm{cap}^{\mathrm{inv}}_n  &= \mathrm{cap}^{\mathrm{inv}}_{\pa(n)}
      + U^{\mathrm{inv}}\, b^{\mathrm{inv}}_n. \label{eq:cap-inv}
\end{align}
The generator, by contrast, is a single active unit whose rating is set by the catalogue
state,
\begin{align}
  \mathrm{cap}^{\mathrm{ge}}_n &= \sum_{s\in\mathcal{S}_{\mathrm{ge}}}
      K^{\mathrm{ge}}_s\, z^{\mathrm{ge}}_{n,s},
  \qquad \sum_{s\in\mathcal{S}_{\mathrm{ge}}} z^{\mathrm{ge}}_{n,s}=1,
  \label{eq:cap-ge}\\
  \mathrm{cap}^{\mathrm{ge}}_n &\geq \mathrm{cap}^{\mathrm{ge}}_{\pa(n)},
  \label{eq:ge-monotone}
\end{align}
where~\eqref{eq:ge-monotone} restricts the generator to upgrades (its rating never
decreases along the tree). The install and retire indicators are linked to the state
transition by
\begin{equation}
  i^{\mathrm{ge}}_{n,s} \geq z^{\mathrm{ge}}_{n,s}-z^{\mathrm{ge}}_{\pa(n),s},
  \qquad
  r^{\mathrm{ge}}_{n,s} \leq z^{\mathrm{ge}}_{\pa(n),s},
  \qquad
  r^{\mathrm{ge}}_{n,s} \leq 1-z^{\mathrm{ge}}_{n,s}.
  \label{eq:ge-transition}
\end{equation}
The installation cost in the objective drives $i^{\mathrm{ge}}$ down to the actual
change, while the resale credit drives $r^{\mathrm{ge}}$ up to its bound, which
equals the actual retirement---so neither can be gamed. The parent state of the root
encodes the initial condition, $\mathrm{cap}^{g}_{\pa(n_0)}=C^{g}_0$ for the modular
technologies and $z^{\mathrm{ge}}_{\pa(n_0),s_0}=1$ for the incumbent generator size $s_0$
(all zero, and $s_0$ the smallest catalogue size, for a greenfield project;
Section~\ref{sec:brownfield}). Because every scenario that shares a common history passes
through the same node, and the investment variables are indexed by node rather than by
scenario, non-anticipativity is enforced implicitly by these recursions: two scenarios
cannot be assigned different capacities before the stage at which they diverge. No
additional coupling constraints are required.

\subsection{Brownfield expansion and generator replacement}
\label{sec:brownfield}

The formulation covers both greenfield projects and---the common case in
practice---\emph{brownfield} projects that expand an already commissioned microgrid. The
two differ only in the root's initial condition. A greenfield project starts empty,
$C^{g}_0=0$; a brownfield project seeds the root's parent with the as-built fleet,
$\mathrm{cap}^{g}_{\pa(n_0)}=C^{g}_0>0$ for the modular technologies and
$z^{\mathrm{ge}}_{\pa(n_0),s_0}=1$ for the incumbent generator. The same expansion
machinery then applies unchanged.

Two economic facts follow automatically. First, the existing fleet is a \emph{sunk cost}:
the capital term $K_n$ (Section~\ref{sec:objective}) prices only the modules \emph{added},
$b^{g}_n$, and the generator \emph{installed}, $i^{\mathrm{ge}}_{n,s}$, so pre-existing
capacity incurs no capital charge. Second, the existing fleet still operates and is
maintained: the operation-and-maintenance term $O_n$ and every dispatch constraint use the
cumulative $\mathrm{cap}^{g}_n$, which includes $C^{g}_0$.

Because the modular technologies only ever add modules, their capacity is monotone
non-decreasing and no retirement is modelled. The generator is different: it is a single
lumpy unit, so expansion means \emph{replacing} the incumbent with a larger catalogue size
rather than running several units at part load. The state
formulation~\eqref{eq:cap-ge}--\eqref{eq:ge-transition} captures exactly this---an upgrade
installs a new size (paying $\sum_s C^{\mathrm{inv}}_{\mathrm{ge},n}K^{\mathrm{ge}}_s
i^{\mathrm{ge}}_{n,s}$) and retires the incumbent for a resale or redeployment credit
($\sum_s V^{\mathrm{ge}}_s r^{\mathrm{ge}}_{n,s}$), with the credit depreciated by the
retired unit's vintage. Restricting the generator to upgrades keeps its rating monotone,
matching the modular technologies.

None of this affects the certificate of Section~\ref{sec:certify}: the capacity lattice is
merely offset so that its lower corner is the existing fleet $C^{g}_0$, and the sunk cost
cancels in both $z_A^\star$ and $z_B^\star$, leaving the price of the heuristic dispatch
$\Delta_{\mathrm{heur}}$ unchanged in structure.

\subsection{Power balance}
\label{sec:balance}

In each operating time step the instantaneous supply must equal the served demand. With
the photovoltaic contribution given by the available yield net of curtailment, the
balance reads
\begin{equation}
  Y_{n,t,h}\,\mathrm{cap}^{\mathrm{pv}}_n - q_{n,t,h}
  + p^{\mathrm{ge}}_{n,t,h}
  + p^{\mathrm{d}}_{n,t,h} - p^{\mathrm{c}}_{n,t,h}
  = D_{n,t,h} - \ell_{n,t,h},
  \qquad \forall (n,t,h).
  \label{eq:balance}
\end{equation}
The product $Y_{n,t,h}\,\mathrm{cap}^{\mathrm{pv}}_n$ couples a parameter to a decision
variable and is therefore linear. Battery discharging adds to supply while charging adds
to the demand side; unmet demand $\ell_{n,t,h}$ relaxes the balance at the cost of the
value of lost load.

\subsection{Generator operation}
\label{sec:generator}

The generator may produce power only when committed, must respect its installed rating,
and---when running---must operate above its minimum stable loading. With $M$ a
sufficiently large constant (conveniently the largest catalogue rating),
\begin{align}
  p^{\mathrm{ge}}_{n,t,h} &\leq \mathrm{cap}^{\mathrm{ge}}_n, \label{eq:ge-cap}\\
  p^{\mathrm{ge}}_{n,t,h} &\leq M\, y_{n,t,h}, \label{eq:ge-on}\\
  p^{\mathrm{ge}}_{n,t,h} &\geq F^{\mathrm{ge}}
      \Big(\sum_{s\in\mathcal{S}_{\mathrm{ge}}}K^{\mathrm{ge}}_s\,
      z^{\mathrm{ge}}_{n,s}\Big) - M\,(1-y_{n,t,h}). \label{eq:ge-min}
\end{align}
Constraint~\eqref{eq:ge-on} forces the output to zero when the unit is off; together with
\eqref{eq:ge-min} it imposes the minimum loading $F^{\mathrm{ge}}$ of the single active
generator when the unit is on, and is vacuous when it is off. Because the generator is one
unit, the minimum loading applies to that unit alone---this is precisely the efficiency a
single, rightly-sized generator provides over several modules idling at part load. The
right-hand side reuses the state variables $z^{\mathrm{ge}}_{n,s}$ of the capacity
definition, and since these are per-node design variables (not per-time-step), no new
per-time-step integer or bilinear term is introduced.

\subsection{Battery storage dynamics}
\label{sec:battery}

The battery energy content evolves from one time step to the next according to a charge
balance that accounts for conversion losses and self-discharge,
\begin{equation}
  e_{n,t,h+1} = (1-A_{n,t,h})\,e_{n,t,h}
  + H^{\mathrm{c}}\,p^{\mathrm{c}}_{n,t,h}
  - \frac{p^{\mathrm{d}}_{n,t,h}}{H^{\mathrm{d}}},
  \qquad \forall (n,t,h).
  \label{eq:soc-dynamics}
\end{equation}
Each representative day is treated as a self-contained cycle, so that the energy content
returns to its initial level at the end of the day,
\begin{equation}
  e_{n,t,0} = e_{n,t,24},
  \qquad \forall (n,t).
  \label{eq:soc-cyclic}
\end{equation}
This cyclic closure removes the need to couple consecutive representative days and keeps
each day's operating sub-problem independent given the installed capacities. The stored
energy is bounded below by the protective discharge trip and above by the
temperature-derated usable capacity,
\begin{equation}
  \underline{E}\,\mathrm{cap}^{\mathrm{batt}}_n
  \;\leq\; e_{n,t,h} \;\leq\;
  \overline{E}\, F^{e}_{n,t,h}\,\mathrm{cap}^{\mathrm{batt}}_n,
  \qquad \forall (n,t,h).
  \label{eq:soc-bounds}
\end{equation}
Charging and discharging powers are limited by the inverter rating and, through it, by
the battery's admissible charge rate,
\begin{equation}
  p^{\mathrm{c}}_{n,t,h} \leq \mathrm{cap}^{\mathrm{inv}}_n,
  \qquad
  p^{\mathrm{d}}_{n,t,h} \leq \mathrm{cap}^{\mathrm{inv}}_n,
  \qquad \forall (n,t,h).
  \label{eq:inv-bounds}
\end{equation}

\subsection{Rule-based generation-balance control}
\label{sec:dispatch}

The operating policy of the microgrid is not cost minimisation but the rule-based
generation-balance controller used in off-grid practice. Its priority order is
straightforward: photovoltaic power serves the load first; any surplus charges the
battery and is curtailed only once the battery is full; remaining demand is met from the
battery; and the generator is started only when the battery alone cannot carry the load.
A look-ahead rule refines this priority at night: the battery is allowed to serve the
load only if it holds enough energy to last until sunrise, and otherwise the generator is
started and used both to serve the load and to recharge the battery up to---but no
further than---the level required for the night.

The central modelling observation is that this controller need not be expressed through
a cascade of logical conditions. The priority order is reproduced by the combination of
the power balance~\eqref{eq:balance}, the generator commitment
constraints~\eqref{eq:ge-cap}--\eqref{eq:ge-min} and the operating economics, because the
objective penalises generator fuel and curtailment; under that objective the optimiser
spontaneously prefers solar, then battery, then generator. The only genuinely additional
behaviour is the look-ahead night reserve, which we impose as a deterministic
time-varying floor on the stored energy,
\begin{equation}
  e_{n,t,h} \;\geq\; R_{n,t,h},
  \qquad \forall (n,t,h).
  \label{eq:night-reserve}
\end{equation}
The reserve $R_{n,t,h}$ is precomputed from the forecast demand between time step $h$ and
the following sunrise, clipped to the installed energy capacity. Constraint
\eqref{eq:night-reserve} forces the controller to keep---or to recharge, using the
generator if necessary---enough energy to carry the community through the night, exactly
as the look-ahead rule prescribes, and it does so without introducing any binary
variable. The generation-balance controller is therefore encoded with a single binary
per time step, namely the generator commitment $y_{n,t,h}$.

Two encodings of the operating layer are retained. In the \emph{baseline} encoding the
constraints \eqref{eq:balance}--\eqref{eq:inv-bounds} are used together with the cost
objective alone; the dispatch is then cost-optimal within the feasible operating
envelope and serves as a tractable reference. In the \emph{rule-faithful} encoding the
night-reserve floor~\eqref{eq:night-reserve} is added, so that the dispatch reproduces the
field controller. The rule-faithful encoding is the model of record; the baseline is used
for validation.

\subsection{Curtailment and unmet demand}
\label{sec:slacks}

Curtailment cannot exceed the available photovoltaic power, and unmet demand cannot
exceed the demand itself,
\begin{equation}
  0 \leq q_{n,t,h} \leq Y_{n,t,h}\,\mathrm{cap}^{\mathrm{pv}}_n,
  \qquad
  0 \leq \ell_{n,t,h} \leq D_{n,t,h},
  \qquad \forall (n,t,h).
  \label{eq:slacks}
\end{equation}

\section{Objective}
\label{sec:objective}

The planner minimises the expected discounted total cost over the scenario tree. The
cost incurred at a node comprises the capital cost of the capacity added there, the
annual operation-and-maintenance cost of the cumulative fleet, and the operating cost of
the representative year---fuel, battery degradation and the penalty on unmet demand---
scaled by the number of operating years the node represents. Writing the incremental
capital cost of node $n$ as
\begin{equation}
  K_n \;=\;
  C^{\mathrm{inv}}_{\mathrm{pv},n}\,U^{\mathrm{pv}} b^{\mathrm{pv}}_n
  + C^{\mathrm{inv}}_{\mathrm{batt},n}\,U^{\mathrm{batt}} b^{\mathrm{batt}}_n
  + C^{\mathrm{inv}}_{\mathrm{inv},n}\,U^{\mathrm{inv}} b^{\mathrm{inv}}_n
  + \sum_{s} C^{\mathrm{inv}}_{\mathrm{ge},n}K^{\mathrm{ge}}_s i^{\mathrm{ge}}_{n,s}
  - \sum_{s} V^{\mathrm{ge}}_s\, r^{\mathrm{ge}}_{n,s},
\end{equation}
in which the modular technologies are charged per added module, a generator upgrade pays
for the newly installed size through $i^{\mathrm{ge}}$ and recovers a resale/transfer
credit on the retired incumbent through $r^{\mathrm{ge}}$, and pre-existing capacity
carries no charge (it is sunk; Section~\ref{sec:brownfield}),
the annual maintenance cost of the cumulative fleet as
\begin{equation}
  O_n \;=\; \sum_{g\in\mathcal{G}} O_g\, C^{\mathrm{inv}}_{g,n}\,\mathrm{cap}^{g}_n,
\end{equation}
and the operating cost of one representative year as
\begin{equation}
  G_n \;=\; \sum_{t\in\mathcal{T}_n} W_{n,t} \sum_{h\in\mathcal{H}}
  \Big[
    C^{\mathrm{fuel}}_n\big(F_0\,y_{n,t,h} + F_1\,p^{\mathrm{ge}}_{n,t,h}\big)
    + C^{\mathrm{deg}}\,p^{\mathrm{d}}_{n,t,h}
    + V^{\mathrm{lol}}\,\ell_{n,t,h}
  \Big],
\end{equation}
the objective is
\begin{equation}
  \min \quad
  \sum_{n\in\mathcal{N}} P_n\,B_n
  \big( K_n + \Delta_n\,O_n + \Delta_n\,G_n \big).
  \label{eq:objective}
\end{equation}
The path probabilities $P_n$ weight each node by the likelihood of its history, so that
\eqref{eq:objective} is the expectation of the discounted life-cycle cost over the
scenario tree---a risk-neutral criterion. Periodic battery replacement and end-of-horizon
salvage are valued in post-processing through the standard net-present-cost accounting and
reported alongside the levelised cost of electricity, defined as the annualised net
present cost divided by the energy served.

\section{Uncertainty space, scenario generation and reduction}
\label{sec:scenarios}

\subsection{The four families of uncertainty}

The uncertain future is described by four families of random factors, each evolving along
the planning horizon.
\begin{description}[leftmargin=2.4cm,style=nextline]
  \item[Demand] The community load is generated bottom-up from appliance ownership and
  usage. Three random trajectories are sampled: the number of connections per customer
  type, the appliance stock per customer type (the name, count and rated power of each
  appliance), and the usage patterns (total functioning time, functioning time per
  switch-on event, time-of-use windows and the probability of occasional use). These feed
  a stochastic load-profile generator calibrated on household surveys and meter readings.
  \item[Renewable resource] The solar resource follows one of the shared
  socio-economic pathways SSP1-2.6, SSP2-4.5, SSP3-7.0 and SSP5-8.5, each of which yields
  a distinct trajectory of irradiance and ambient temperature and hence of the specific
  yield $Y$.
  \item[Economic conditions] The diesel price and the investment-cost trajectories of the
  photovoltaic array, the battery, the generator and the inverter are sampled, capturing
  fuel-market volatility and technology learning.
  \item[Policy] A minimum solar-penetration target may be imposed, sampled as a discrete
  policy factor.
\end{description}

\subsection{Monte-Carlo sampling, scenario reduction and tree construction}

A large Monte-Carlo ensemble of complete scenario paths is drawn by sampling each family
of uncertainty along the horizon. Because the resulting ensemble is far too large to
optimise over directly, it is condensed in two complementary ways.

First, the ensemble is reduced \emph{across scenarios}. At each stage the sampled
outcomes are clustered---by a medoid-based reduction such as $k$-medoids or fast-forward
selection on a feature representation of the paths---to a small set of representative
outcomes, each carrying the probability mass of its cluster. Nesting the stage-wise
representative outcomes produces a branching scenario tree: a node at one stage gives
rise to several children at the next, and the path probability $P_n$ is the product of
the conditional branch probabilities along the path. A two-stage program is recovered as
the special case of a tree with a single decision node at the first stage.

Second, the ensemble is reduced \emph{in time}. Within each node, the full operating year
is compressed into a handful of weighted representative days (typically four to twelve) by
clustering the daily demand and resource shapes, the medoid of each cluster becoming a
representative day $t$ with weight $W_{n,t}$ equal to the size of its cluster. This
time-domain reduction, combined with the cyclic closure~\eqref{eq:soc-cyclic}, is what
keeps the operating sub-problems small enough that the whole tree can be optimised at
once. The error introduced by both reductions, relative to the full ensemble, is
quantified and reported so that the compression remains transparent.

\section{Deterministic equivalent and solution}
\label{sec:solution}

Collecting the node-indexed investment variables, the time-step-indexed operating
variables, the constraints of Section~\ref{sec:constraints} over all nodes and time
steps, and the objective~\eqref{eq:objective} yields a single mixed-integer linear
program---the deterministic equivalent of the multi-stage stochastic program. Its only
integer variables are the modular module counts $b^{\mathrm{pv}}_n$,
$b^{\mathrm{batt}}_n$, $b^{\mathrm{inv}}_n$, the generator state
$z^{\mathrm{ge}}_{n,s}$, and one generator-commitment binary per operating time step. The number of commitment binaries is
the product of the number of tree nodes, the number of representative days per node and
the twenty-four hours of the day; with a modest number of representative days this remains
well within the reach of contemporary mixed-integer solvers. The non-anticipativity of the
investment policy is, as noted, embedded in the tree indexing and requires no explicit
constraints. The model is implemented in an algebraic modelling layer over labelled
multidimensional arrays and solved by a standard branch-and-cut solver.

\section{Certified-optimal sizing under the rule-based dispatch policy}
\label{sec:certify}

The deterministic equivalent of Section~\ref{sec:solution} returns a globally optimal
sizing \emph{under the assumption that operation is cost-optimal}. In off-grid practice,
however, the controller that will actually be deployed is the rule-based
generation-balance policy of Section~\ref{sec:dispatch}, which is in general
sub-optimal. The sizing that is optimal for the deployed controller therefore differs
from the sizing that is optimal for cost-minimising dispatch, and it is the former that a
planner ultimately needs. This section shows that one can certify global optimality of
the sizing \emph{for the fixed rule-based controller}, without encoding the controller as
constraints, by using the cost-optimal dispatch problem as a rigorous lower bound and the
forward simulation of the controller as an upper bound.

\subsection{Two operating cost functions}

Let $x$ denote the complete vector of investment decisions---equivalently the cumulative
capacities $\{\mathrm{cap}^{g}_n\}_{g,n}$---and let $\mathcal{X}$ be the finite discrete
lattice of admissible investments (integer PV, battery and inverter module counts within
bounds---floored at the existing fleet $C^{g}_0$---and the single generator's catalogue
size). For a fixed $x$ and a node $n$, write
$\mathcal{F}_n(x)$ for the operating feasible set defined by the power
balance~\eqref{eq:balance}, the generator
constraints~\eqref{eq:ge-cap}--\eqref{eq:ge-min}, the storage
dynamics~\eqref{eq:soc-dynamics}--\eqref{eq:inv-bounds} and the
slacks~\eqref{eq:slacks}, and let $c_n(\cdot)$ be the per-year operating cost $G_n$. Define
\begin{align}
  C^{\mathrm{opt}}_n(x) &= \min_{u\,\in\,\mathcal{F}_n(x)} c_n(u)
      && \text{(cost-optimal dispatch),}\\
  C^{\mathrm{rule}}_n(x) &= c_n\big(u^{\mathrm{rule}}_n(x)\big)
      && \text{(rule-based dispatch),}
\end{align}
where $u^{\mathrm{rule}}_n(x)$ is the trajectory produced by the generation-balance
controller. With $K_n,O_n$ as in Section~\ref{sec:objective}, the two sizing objectives are
\begin{equation}
  Z_\bullet(x) = \sum_{n\in\mathcal{N}} P_n\,B_n
  \big( K_n(x) + \Delta_n O_n(x) + \Delta_n\,C^{\bullet}_n(x) \big),
  \qquad \bullet\in\{\mathrm{opt},\mathrm{rule}\},
\end{equation}
with global optima $z_A^\star=\min_{x\in\mathcal{X}} Z_{\mathrm{opt}}(x)$ (the
deterministic equivalent of Section~\ref{sec:solution}) and
$z_B^\star=\min_{x\in\mathcal{X}} Z_{\mathrm{rule}}(x)$ (the quantity of interest).

\subsection{The dispatch-relaxation lower bound}

\begin{assumption}[Policy feasibility]
\label{ass:feas}
For every $x\in\mathcal{X}$ and node $n$, the rule-based trajectory lies in the operating
feasible set, $u^{\mathrm{rule}}_n(x)\in\mathcal{F}_n(x)$.
\end{assumption}

Assumption~\ref{ass:feas} merely states that the controller respects the same physics as
the cost-optimal model. It holds for the generation-balance controller of
Section~\ref{sec:dispatch}: the controller satisfies the power balance with the unmet-load
slack, keeps the state of charge within the trip limits~\eqref{eq:soc-bounds}, runs the
generator above its minimum loading or not at all~\eqref{eq:ge-min}, respects the inverter
limits~\eqref{eq:inv-bounds}, and never charges and discharges simultaneously; its
trajectory is thus one feasible point of $\mathcal{F}_n(x)$, not necessarily the
cost-minimising one.

\begin{proposition}[Dispatch-relaxation bound]
\label{prop:bound}
Under Assumption~\ref{ass:feas}, $C^{\mathrm{opt}}_n(x)\le C^{\mathrm{rule}}_n(x)$ for all
$x\in\mathcal{X}$ and all $n$. Consequently $Z_{\mathrm{opt}}(x)\le Z_{\mathrm{rule}}(x)$
for every $x$, and $z_A^\star\le z_B^\star$. Moreover, for any box
$B\subseteq\operatorname{conv}(\mathcal{X})$, the continuous-capacity relaxation
$\underline{Z}(B)=\min_{x\in B}Z_{\mathrm{opt}}(x)$ is a valid lower bound on
$\min_{x\in\mathcal{X}\cap B} Z_{\mathrm{rule}}(x)$.
\end{proposition}

\begin{proof}
By Assumption~\ref{ass:feas}, $u^{\mathrm{rule}}_n(x)$ is feasible for the minimisation
defining $C^{\mathrm{opt}}_n(x)$, hence
$C^{\mathrm{opt}}_n(x)=\min_{u\in\mathcal{F}_n(x)}c_n(u)\le
c_n(u^{\mathrm{rule}}_n(x))=C^{\mathrm{rule}}_n(x)$. The investment and maintenance terms
$K_n,O_n$ are identical in $Z_{\mathrm{opt}}$ and $Z_{\mathrm{rule}}$, and the weights
$P_n B_n\Delta_n\ge0$; summing preserves the inequality, giving
$Z_{\mathrm{opt}}(x)\le Z_{\mathrm{rule}}(x)$ pointwise and $z_A^\star\le z_B^\star$ after
minimisation. For the relaxation, any $x\in\mathcal{X}\cap B$ satisfies
$Z_{\mathrm{rule}}(x)\ge Z_{\mathrm{opt}}(x)\ge\min_{x'\in B}Z_{\mathrm{opt}}(x')
=\underline{Z}(B)$, since $\mathcal{X}\cap B\subseteq B$ and integrality is relaxed.
\end{proof}

The bound needs no monotonicity argument: relaxing integrality (a superset of capacities)
and replacing the rule-based cost by the cost-optimal cost (a smaller value) both can only
\emph{decrease} the objective, so $\underline{Z}(B)$ underestimates the rule-based optimum
over the box. The lower-bounding problem $\underline{Z}(B)$ is exactly the
deterministic-equivalent program of Section~\ref{sec:solution} restricted to $x\in B$;
relaxing the generator-commitment binary in it yields a still-valid but cheaper linear
lower bound.

\subsection{Branch-and-simulate}

Proposition~\ref{prop:bound} pairs a computable lower bound (solve the cost-optimal
relaxation on a box) with a computable upper bound (forward-simulate the controller at an
integer point), which is exactly what a branch-and-bound certificate requires.

\begin{corollary}[Certificate]
\label{cor:cert}
For any incumbent $\hat{x}\in\mathcal{X}$ and any partition $\{B_k\}$ of
$\operatorname{conv}(\mathcal{X})$,
$\min_k \underline{Z}(B_k)\le z_B^\star\le Z_{\mathrm{rule}}(\hat{x})$. If
$Z_{\mathrm{rule}}(\hat{x})-\min_k\underline{Z}(B_k)\le\varepsilon$, then $\hat{x}$ is an
$\varepsilon$-optimal sizing for the rule-based controller.
\end{corollary}

This gives the following procedure, which we call \emph{branch-and-simulate}.
\begin{enumerate}[leftmargin=2em,itemsep=1pt]
  \item Initialise the incumbent value $\bar{z}\leftarrow+\infty$ and a queue of open
  boxes with the root $B_0=\operatorname{conv}(\mathcal{X})$.
  \item Pop a box $B$. Compute the lower bound $\underline{Z}(B)$ by solving the
  cost-optimal relaxation (Section~\ref{sec:solution}) restricted to $x\in B$.
  \item If $\underline{Z}(B)\ge\bar{z}-\varepsilon$, \emph{prune} $B$
  (Proposition~\ref{prop:bound}).
  \item Otherwise round the relaxed solution to an integer point $\hat{x}\in
  \mathcal{X}\cap B$, evaluate the upper bound $Z_{\mathrm{rule}}(\hat{x})$ by
  \emph{forward-simulating} the controller on every node and scenario, and update the
  incumbent if it improves.
  \item Branch: split $B$ along an integer capacity coordinate and return the children to
  the queue.
  \item Terminate when the queue is empty or
  $\bar{z}-\min_{\text{open }B}\underline{Z}(B)\le\varepsilon$; return the incumbent and
  its certified gap.
\end{enumerate}
Because $\mathcal{X}$ is finite the branching terminates, and on termination the incumbent
is a certified $\varepsilon$-global optimum of the rule-based sizing problem $z_B^\star$.
Only linear or mixed-integer relaxations (lower bounds) and forward simulations of the
real controller (upper bounds) are ever solved; the controller is never approximated by
constraints, so its behaviour is reproduced \emph{exactly}. The pruning is tight precisely
when the rule-based controller is close to cost-optimal, so that the method is fast in the
regime where the heuristic is good and merely slower---never incorrect---when it is poor.

\subsection{The price of the heuristic dispatch}

\begin{definition}[Price of the heuristic dispatch]
\label{def:price}
$\;\Delta_{\mathrm{heur}} = z_B^\star - z_A^\star \ge 0$, reported in absolute terms and as
the relative gap $\Delta_{\mathrm{heur}}/z_A^\star$.
\end{definition}

By Proposition~\ref{prop:bound} this quantity is non-negative; it measures the avoidable
life-cycle cost attributable to operating the microgrid with the rule-based controller
instead of an ideal cost-optimal dispatcher, evaluated at each method's own optimal
sizing. It is a by-product of the same two solves---$z_A^\star$ is the
deterministic-equivalent optimum and $z_B^\star$ the branch-and-simulate optimum---and is
itself an informative result: a small $\Delta_{\mathrm{heur}}$ justifies the operational
simplicity of the heuristic, whereas a large one quantifies the value of investing in
smarter control.

\subsection{Refining the controller: a parameterised rule class}
\label{sec:refine}

The generation-balance controller of Section~\ref{sec:dispatch} is one point in a much
larger space of deployable, causal operating policies. Rather than commit to a single set
of rules, we embed it in a \emph{parameterised policy class}
$\{\pi_\theta : \theta\in\Theta\}$, where each $\pi_\theta$ is a causal controller that,
at every time step, maps the current battery energy and the information available up to
that step to a feasible dispatch. Four parameters capture the levers that matter most for
off-grid operating cost.
\begin{description}[leftmargin=1.9em,style=nextline]
  \item[Commitment hysteresis $(\tilde E_{\mathrm{on}},\tilde E_{\mathrm{off}})$.] With
  $\underline{E}\le\tilde E_{\mathrm{on}}\le\tilde E_{\mathrm{off}}\le\overline E$, the
  generator starts when the state of charge would fall below
  $\tilde E_{\mathrm{on}}F^{e}\,\mathrm{cap}^{\mathrm{batt}}$ and, once started, keeps
  running until the charge reaches $\tilde E_{\mathrm{off}}F^{e}\,\mathrm{cap}^{\mathrm{batt}}$.
  A wide band suppresses short, inefficient generator cycles.
  \item[{Generator setpoint $T^{\mathrm{ge}}\in[F^{\mathrm{ge}},1]$.}] When the
  generator runs it targets a loading $T^{\mathrm{ge}}$ of its rating; the power in
  excess of the served deficit charges the battery. $T^{\mathrm{ge}}=F^{\mathrm{ge}}$
  recovers pure load-following, while larger values realise cycle-charging that runs the
  unit near its efficient point.
  \item[Reserve multiplier $M^{\mathrm{r}}\ge0$.] The night-reserve floor is scaled,
  $R_{n,t,h}(\theta)=M^{\mathrm{r}}\,\hat R_{n,t,h}$, where $\hat R_{n,t,h}$ is the forecast net
  demand until the next photovoltaic surplus; $M^{\mathrm{r}}$ trades reliability against fuel and
  subsumes the fixed reserve of Section~\ref{sec:dispatch} at a nominal $M^{\mathrm{r}}_0$.
  \item[Look-ahead horizon $L\ge0$.] The number of future steps the policy anticipates
  through a forecast of demand and yield; $L=0$ is purely reactive and increasing $L$
  yields predictive (receding-horizon) control.
\end{description}
Named strategies are special cases: load-following is
$T^{\mathrm{ge}}=F^{\mathrm{ge}}$ with $\tilde E_{\mathrm{on}}=\tilde E_{\mathrm{off}}$;
cycle-charging raises $T^{\mathrm{ge}}$ and widens the band; a combined policy selects
the cheaper of the two each step; and predictive control corresponds to $L>0$. The class
$\Theta$ is constructed so that every $\pi_\theta$ produces trajectories in the feasible
set $\mathcal{F}(x)$---it never charges and discharges simultaneously, respects the
minimum loading $F^{\mathrm{ge}}$ and the state-of-charge trips---so
Assumption~\ref{ass:feas} holds for all $\theta\in\Theta$.

\paragraph{The bound is invariant to the controller.} This last point has a decisive
consequence. By Proposition~\ref{prop:bound}, $C^{\mathrm{opt}}_n(x)\le
C^{\mathrm{rule}}_n(x;\theta)$ for \emph{every} feasible $\theta$, because each $\pi_\theta$
is one feasible dispatch. Hence the lower bound $z_A^\star$ and the entire
branch-and-simulate certificate of Section~\ref{sec:certify} are \emph{independent of the
choice of controller}: refining the rules can only \emph{lower} the upper bound
$z_B(\theta)$, never invalidate the certificate. Controller refinement is therefore
``free'' with respect to correctness---a better rule simply shrinks the certified price of
the heuristic dispatch.

\paragraph{Gap-minimising tuning.} The best controller in the class, for a given sizing
$x$, minimises the expected discounted operating cost,
\begin{equation}
  \theta^\star(x) \;=\; \arg\min_{\theta\in\Theta}\;
  \sum_{n\in\mathcal{N}} P_n\,B_n\,\Delta_n\,C^{\mathrm{rule}}_n(x;\theta),
  \label{eq:tune}
\end{equation}
which, since $C^{\mathrm{rule}}$ is evaluated by forward simulation, is solved by a
derivative-free search over the low-dimensional $\Theta$ (grid, Nelder--Mead or
differential evolution). Crucially, the cost-optimal dispatch acts as a \emph{teacher}:
comparing the optimal trajectory $u^{\mathrm{opt}}_n(x)$ with the trajectory of $\pi_\theta$
exposes the exact time steps at which the rule loses money, which both warm-starts $\theta$
(imitation of the optimal policy) and guides the choice of which parameters to enrich.
Substituting $\theta^\star(x)$ into the upper-bound oracle of branch-and-simulate tightens
$z_B^\star$ and, jointly with the sizing search, co-selects the pair $(x,\theta)$.

\paragraph{Achievable price and its decomposition.} For a policy class $\Theta$ define the
class-residual price
$\Delta_{\mathrm{heur}}(\Theta)=\min_{\theta\in\Theta} z_B(\theta)-z_A^\star\ge0$, which is
non-increasing as $\Theta$ is enriched and tends to zero as the class approaches the
full-foresight cost-optimal policy. It answers, quantitatively and with a certificate, the
practical question \emph{``how close to optimal can a simple deployable rule get?''}
Writing $\Theta_{\mathrm{c}}$ for the causal (finite-foresight) sub-class, the price splits
as
\begin{equation}
  \Delta_{\mathrm{heur}}(\theta)
  = \underbrace{\big[z_B(\theta)-\min_{\vartheta\in\Theta_{\mathrm{c}}}z_B(\vartheta)\big]}
      _{\text{sub-optimality of the chosen rule}}
  + \underbrace{\big[\min_{\vartheta\in\Theta_{\mathrm{c}}}z_B(\vartheta)-z_A^\star\big]}
      _{\text{value of foresight}} ,
\end{equation}
separating the cost of using a simple rule (attacked by the hysteresis, setpoint and
reserve parameters) from the cost of imperfect information (attacked by the look-ahead
horizon $L$). The decomposition tells the planner whether to invest in smarter logic or in
better forecasting.

\subsection{Relation to prior work}

The literature on off-grid microgrid sizing separates into two families that this
construction bridges. The first uses rule-based controllers---HOMER's load-following,
cycle-charging, generator-order and combined strategies, and the genetic-algorithm tool
iHOGA/MHOGA of Dufo-L\'opez and Bernal-Agust\'in---but searches the capacity space by
enumeration or metaheuristics and reports \emph{no} optimality certificate or gap. The
second uses exact mathematical programming---e.g.\ the MILP planning model of Wu \emph{et
al.}\ (IEEE Trans.\ Smart Grid 12(5), 2021), the MicroGridsPy capacity-expansion framework
of Stevanato \emph{et al.}, and the multi-stage stochastic integer template of Ahmed,
King \& Parija (J.\ Glob.\ Optim.\ 26, 2003)---which certifies optimality but only for
\emph{cost-optimal} dispatch. The closest hybrids confirm rather than pre-empt the gap:
the inner--outer co-optimisation of Lyu \emph{et al.}\ (Applied Energy 363, 2024) keeps an
optimiser in the lower level and explicitly deems sizing by rule-based strategies
``risky''; the progressive-hedging scheme of Zolan/Wales \emph{et al.}\ (Optimization and
Engineering 26, 2025) yields an optimality gap but bounds the cost-optimal-dispatch MILP
itself; and the rule-based controller tuning of Nespoli \& Medici (arXiv:2511.21619, 2025)
uses a metaheuristic with no certificate and explicitly dismisses a nested
master/inner-optimisation as too costly. To our knowledge, no prior work certifies global
optimality of the sizing \emph{for a fixed rule-based controller} by using cost-optimal
dispatch as a guaranteed lower bound, nor quantifies the price of the heuristic dispatch as
a certified gap. Two qualifications are in order: the certificate rests on
Assumption~\ref{ass:feas}, which must be verified for any controller to which the method is
applied, and the extension of the bound to a genuinely adaptive multi-stage recourse---in
which a fixed (non-optimising) policy would remove the integer recourse that obstructs
stochastic dual dynamic integer programming---is a promising but as yet unproven direction.

\end{document}
